\chapter{问卷调查结果及分析}


\section{描述性分析}

本次调查要求数学专业师范生对自身的教学技能掌握状况进行评判和评估，得出的结论为师范生自我报告的数学教学技能水平。


\subsection{数学专业师范生教学技能掌握现状}

将调查问卷第二部分所得的结果进行描述性统计（见表~\ref{table:数学教学技能的描述性统计} ），在此基础上对数学师范生的教学技能水平进行全面分析。

\begin{table}[htbp]
  \centering
  \caption{数学教学技能的描述性统计}
  \label{table:数学教学技能的描述性统计}
  \begin{tblr}{
    width = 0.7\textwidth,
    colspec = {X[2,c]X[2,c]X[1,c]},
    hline{1,2,Z} = {solid},
  }
    项目     & 平均值  & 标准差  \\
    教学设计技能 & 3.5  & 0.55 \\
    数学说课技能 & 3.32 & 0.62 \\
    课堂教学技能 & 3.53 & 0.45 \\
    教学评价技能 & 3.31 & 0.65 \\
    教学研究技能 & 3.34 & 0.66 \\
    数学教学技能 & 3.4  & 0.45
  \end{tblr}
\end{table}


表~\ref{table:数学教学技能的描述性统计} 表明，教学设计和课堂教学技能得分较高，其中课堂教学技能得分最高，为3.53；其他教学技能得分明显低于前两项技能，其中教学评价技能得分最低。这说明数学师范生的教学设计和课堂教学技能水平整体高于数学说课、教学评价和教学研究技能水平。

\begin{table}[htbp]
  \centering
  \caption{课堂教学技能的描述性统计}
  \label{table:课堂教学技能的描述性统计}
  \begin{tblr}{
    width = 0.7\textwidth,
    colspec = {X[2,c]X[2,c]X[1,c]},
    hline{1,2,Z} = {solid},
  }
    项目     & 平均值  & 标准差  \\
    语言技能   & 3.45 & 0.61 \\
    板书技能   & 3.51 & 0.7  \\
    讲解技能   & 3.61 & 0.72 \\
    提问技能   & 3.49 & 0.57 \\
    导入技能   & 3.52 & 0.55 \\
    结束技能   & 3.61 & 0.56 \\
    媒体运用技能 & 3.53 & 0.59 \\
  \end{tblr}
\end{table}


表~\ref{table:课堂教学技能的描述性统计} 表明，在所有的课堂教学技能中，讲解和结束技能得分最高，为3.61；其次是媒体运用技能、导入技能和板书技能，得分均在3.5以上；而提问技能和语言技能得分偏低。这说明师范生的提问技能和语言技能还有待加强。


\subsection{数学教学技能训练现状分析}

问卷第三部分调查了学校各项教学技能的培养情况，以及数学专业师范生对各种培养方式的作用的看法，具体数据见表7和表8。
表7统计了数学师范生本科期间各项教学技能训练次数，具体结果见表7：

\begin{table}[htbp]
  \centering
  \caption{教学技能训练次数统计}
  \label{table:教学技能训练次数统计}
  \begin{tblr}{
    width = \textwidth,
    colspec = {X[14em,c]X[1,c]X[1,c]X[1,c]X[1,c]},
    hline{1,2,Z} = {solid},
  }
    \diagbox{题目}{选项\hspace*{4.5em}}          & 0次   & 0-5次 & 5-10次 & 10次以上 \\
    本科期间进行教学设计次数   & 0\%  & 31\% & 36\%  & 33\%  \\
    本科期间进行说课次数     & 17\% & 68\% & 11\%  & 5\%   \\
    本科期间进行讲课次数     & 0\%  & 38\% & 40\%  & 22\%  \\
    本科期间参与评课次数     & 9\%  & 38\% & 28\%  & 25\%  \\
    本科期间参与教育教学研究次数 & 20\% & 61\% & 14\%  & 6\%   \\
  \end{tblr}
\end{table}

表~\ref{table:教学技能训练次数统计} 表明：本科期间所有数学师范生都进行了教学设计训练和讲课训练，且训练次数相对较多。教学设计次数和讲课次数在5次以上的人数占比均超过60\%，其中教学设计10次以上的人数占比33\%，在所有数学教学技能中排第一位。这项数据说明，学校对这两项教学技能的训练较为充分，尤其是教学设计技能的训练次数最多，参与学生也最多。

少部分师范生没有进行过说课训练，甚至有师范生没有参与过评课和教育教学研究。大部分师范生的说课次数和教育教学研究次数都集中在0-5次，只有不到10\%的数学师范生说课次数和教育教学研究次数达到5次以上。数学师范生的评课次数比例相对较分散，评课次数在0-5次的人数偏多，评课5-10次和10次以上的人数差距很小。

表~\ref{table:师范生培养活动对各项教学技能的培养效果} 统计了学校各项师范生培养活动对各项教学技能的培养效果。

\begin{table}[htbp]
  \centering
  \caption{师范生培养活动对各项教学技能的培养效果}
  \label{table:师范生培养活动对各项教学技能的培养效果}
  \begin{tblr}{
    width = 1.1\textwidth,
    % colspec = {X[3.7,c]X[1,c]X[1,c]X[1,c]X[1,c]X[1,c]},
    % colspec = {X[3.8,c]X[6em,c]X[6em,c]X[6em,c]X[6em,c]X[6em,c]},
    colspec = {X[3.8,c]X[4em,c]X[4em,c]X[4em,c]X[4em,c]X[4em,c]},
    hline{1,2,Z} = {solid},
    stretch = 0.9,
    row{1} = {m}
  }
    \diagbox{
      \raisebox{2.2em}{题目}
      % \rule{0pt}{4em} 题目
    }{选项\hspace*{4em}} & 教学设计技能 & 数学说课技能 & 课堂教学技能 & 教学评价技能 & 教学研究技能 \\
    你认为教学比赛和师范生技能比赛培养了哪方面教学技能 & 82\%   & 66\%   & 86\%   & 46\%   & 36\%   \\
    {你认为教育见习和教育实习\\培养了哪方面教学技能}    & 75\%   & 51\%   & 91\%   & 61\%   & 52\%   \\
    {你认为微格教学培养了哪方面\\教学技能}         & 73\%   & 46\%   & 79\%   & 64\%   & 35\%   \\
    {你认为理论学习培养了哪方面\\教学技能}         & 50\%   & 37\%   & 47\%   & 53\%   & 59\%   \\
    {你认为观摩学习培养了哪方面\\教学技能}         & 60\%   & 49\%   & 77\%   & 63\%   & 34\% 
  \end{tblr}
\end{table}

从横向上分析表~\ref{table:师范生培养活动对各项教学技能的培养效果} ，教学比赛和师范生技能比赛主要培养了师范生的教学设计、课堂教学和数学说课技能；教育见习和实习对各项数学教学技能的提高都有一定的促进作用，对课堂教学技能的培养效果最好；微格教学主要培养了师范生的教学设计、课堂教学和教学评价技能；理论学习（教育学、心理学、数学教学论等）对各项教学技能的提升作用相对较小，但对教学研究技能的培养作用较大；观摩学习主要培养了课堂教学技能，其次是教学设计技能和教学评价技能。

从纵向上分析表~\ref{table:师范生培养活动对各项教学技能的培养效果} ，学校的教学比赛、教育见习和实习、微格教学、理论学习以及观摩学习均能够有效的培养师范生的教学设计技能和课堂教学技能；培养数学说课技能最有效的活动是教学比赛，其次是教育见习、实习和观摩学习；各项活动对教学评价和教学研究技能的培养作用都相对较低。

从整体上看，数学师范生普遍认为学校的教学比赛、教育见习和实习、微格教学、理论学习以及观摩学习能够有效的培养他们的各项教学技能，其中教学设计和课堂教学技能的训练效果最明显。



\section{差异性分析}

本研究主要采用独立样本T检验和单因素方差分析对问卷数据进行差异性分析。


\subsection{年级差异}

本研究比较了师范生的数学教学技能在年级上的差异，统计结果如表~\ref{table:各个维度在年级上的差异分析结果} 所示。

\begin{table}
  \centering
  \caption{各个维度在年级上的差异分析结果}
  \label{table:各个维度在年级上的差异分析结果}
  \begin{tblr}{
    width  = \textwidth,
    colspec = {
      X[2,c]X[2em,c]X[3em,c]X[3em,c]X[4em,c]X[1,c]X[3em,c]
    },
    hline{1,2,Z} = {solid},
    cell{2}{1} = {r=2}{m},
    cell{4}{1} = {r=2}{m},
    cell{6}{1} = {r=2}{m},
    cell{8}{1} = {r=2}{m},
    cell{10}{1} = {r=2}{m},
    cell{12}{1} = {r=2}{m},
    cell{2}{6,7} = {r=2}{m},
    cell{4}{6,7} = {r=2}{m},
    cell{6}{6,7} = {r=2}{m},
    cell{8}{6,7} = {r=2}{m},
    cell{10}{6,7} = {r=2}{m},
    cell{12}{6,7} = {r=2}{m},
  }
    变量     & 年级 & 个案数 & 平均值  & 标准偏差  & t     & 显著性   \\
    教学设计技能 & 大四 & 106 & 3.61 & 0.5   & 3.413 & 0.001 \\
          & 大三 & 71  & 3.33 & 0.57  &       &       \\
    数学说课技能 & 大四 & 106 & 3.33 & 0.55  & 0.252 & 0.801 \\
          & 大三 & 71  & 3.31 & 0.71  &       &       \\
    课堂教学技能 & 大四 & 106 & 3.58 & 0.42  & 1.682 & 0.094 \\
          & 大三 & 71  & 3.46 & 0.48  &       &       \\
    教学评价技能 & 大四 & 106 & 3.36 & 0.6   & 1.21  & 0.228 \\
          & 大三 & 71  & 3.23 & 0.71  &       &       \\
    教学研究技能 & 大四 & 106 & 3.36 & 0.63  & 0.654 & 0.514 \\
          & 大三 & 71  & 3.3  & 0.7   &       &       \\
    数学教学技能 & 大四 & 106 & 3.45 & 0.405 & 1.76  & 0.08  \\
          & 大三 & 71  & 3.32 & 0.503 &       &  
  \end{tblr}
\end{table}

从表~\ref{table:各个维度在年级上的差异分析结果} 可以看出：（1）在随机抽取的177名数学师范生中，大三和大四的师范生在数学教学技能上的差异不明显，大四的数学教学技能的平均分为3.45，大三的平均分为3.32，大四师范生的整体素质略高于大三师范生。（2）单看每一项教学技能，可以发现教学设计技能在年级上的差异显著性检验为0.001，明显小于0.05，说明不同年级的师范生在教学设计技能上存在差异，且大四明显高于大三。而其他教学技能在年级上不存在显著的统计学差异，因为sig均大于0.05，但大四师范生的各项教学技能平均分均高于大三师范生。


\subsection{性别差异}

本研究对师范生数学教学技能的性别差异进行了对比分析，具体数据如表~\ref{table:各个维度在性别上的差异分析结果} 所示。

\begin{table}
  \centering
  \caption{各个维度在性别上的差异分析结果}
  \label{table:各个维度在性别上的差异分析结果}
  \begin{tblr}{
    width  = \textwidth,
    colspec = {
      X[2,c]X[2em,c]X[3em,c]X[3em,c]X[4em,c]X[1,c]X[3em,c]
    },
    hline{1,2,Z} = {solid},
    cell{2}{1} = {r=2}{m},
    cell{4}{1} = {r=2}{m},
    cell{6}{1} = {r=2}{m},
    cell{8}{1} = {r=2}{m},
    cell{10}{1} = {r=2}{m},
    cell{12}{1} = {r=2}{m},
    cell{2}{6,7} = {r=2}{m},
    cell{4}{6,7} = {r=2}{m},
    cell{6}{6,7} = {r=2}{m},
    cell{8}{6,7} = {r=2}{m},
    cell{10}{6,7} = {r=2}{m},
    cell{12}{6,7} = {r=2}{m},
  }
    变量     & 性别 & 个案数 & 平均值  & 标准偏差 & t     & 显著性   \\
    教学设计技能 & 女  & 141 & 3.55 & 0.52 & 2.372 & 0.019 \\
          & 男  & 36  & 3.31 & 0.62 &       &       \\
    数学说课技能 & 女  & 141 & 3.35 & 0.57 & 1.381 & 0.169 \\
          & 男  & 36  & 3.19 & 0.76 &       &       \\
    课堂教学技能 & 女  & 141 & 3.57 & 0.45 & 2.354 & 0.02  \\
          & 男  & 36  & 3.38 & 0.43 &       &       \\
    教学评价技能 & 女  & 141 & 3.35 & 0.65 & 1.551 & 0.123 \\
          & 男  & 36  & 3.16 & 0.64 &       &       \\
    教学研究技能 & 女  & 141 & 3.35 & 0.67 & 0.398 & 0.691 \\
          & 男  & 36  & 3.3  & 0.61 &       &       \\
    数学教学技能 & 女  & 141 & 3.43 & 0.43 & 1.991 & 0.048 \\
          & 男  & 36  & 3.26 & 0.5  &       &       \\
  \end{tblr}
\end{table}


从表~\ref{table:各个维度在性别上的差异分析结果} 可以看出：（1）在随机抽取的177名数学师范生中，男生和女生在数学教学技能上的差异显著性检验为0.048，小于0.05，说明男生和女生在数学教学技能上存在差异，且女生数学教学技能掌握情况比男生好。（2）单看每一项教学技能，可以发现教学设计技能和课堂教学技能在性别上的差异显著性检验为0.019和0.02，明显小于0.05，说明不同性别的师范生在教学设计技能和课堂教学技能上存在差异，且女生平均分高于男生。（3）而数学说课、教学评价和教学研究技能在性别上不存在显著的统计学差异，因为sig均大于0.05，但女生的各项教学技能平均分均高于男生。


\subsection{普通话等级差异}

本研究比较了普通话等级不同的师范生在数学教学技能上差异，统计结果如表~\ref{table:各个维度在普通话等级上的差异分析结果} 所示。

\begin{table}
  \centering
  \caption{各个维度在普通话等级上的差异分析结果}
  \label{table:各个维度在普通话等级上的差异分析结果}
  \begin{tblr}{
    width  = \textwidth,
    colspec = {
      X[2,c]X[1,c]X[1,c]X[3,c]
    },
    hline{1,2,Z} = {solid},
  }
    变量     & F     & 显著性   & 多重比较 \\
    教学设计技能 & 4.904 & 0.003 &  一乙$>$二甲，一乙$>$二乙*    \\
    数学说课技能 & 3.199 & 0.025 &  一乙$>$二甲，一乙$>$二乙*    \\
    课堂教学技能 & 8.002 & 0     &  一乙$>$二甲，一乙$>$二乙*    \\
    教学评价技能 & 4.954 & 0.003 &  一乙$>$二甲，一乙$>$二乙*    \\
    教学研究技能 & 2.018 & 0.113 & 无显著差异     \\
    数学教学技能 & 6.584 & 0     &  一乙$>$二甲，一乙$>$二乙*  
  \end{tblr}
  *在0.05的显著性水平上平均数差异显著。
\end{table}


从表~\ref{table:各个维度在普通话等级上的差异分析结果} 可以看出：（1）普通话等级不同的师范生在数学教学技能上差异显著性检验为0.000，明显小于0.05，说明普通话等级不同的师范生在数学教学技能上存在显著的统计学差异，且普通话等级为一乙的师范生平均分高于普通话等级为二甲和二乙的师范生。（2）单看每一项教学技能，可以发现计算机等级不同的师范生在教学设计、数学说课、课堂教学和教学评价技能上的sig值均小于0.05，说明普通话等级不同的师范生在这些技能上存在差异，普通话等级为一乙的师范生平均分高于普通话等级为二甲和二乙的师范生。而普通话等级不同的师范生在教学研究技能上不存在显著的差异，因为sig大于0.05。


\subsection{计算机等级差异}

本研究比较了计算机等级不同的师范生在数学教学技能上的差异，统计结果如表~\ref{table:各维度在计算机等级上的差异性分析结果} 所示。

\begin{table}
  \centering
  \caption{各维度在计算机等级上的差异性分析结果}
  \label{table:各维度在计算机等级上的差异性分析结果}
  \begin{tblr}{
    width  = \textwidth,
    colspec = {
      X[2,c]X[1,c]X[1,c]X[3,c]
    },
    hline{1,2,Z} = {solid},
  }
    变量     & F     & 显著性   & 多重比较  \\
    教学设计技能 & 1.418 & 0.245 & 无显著差异 \\
    数学说课技能 & 1.06  & 0.349 & 无显著差异 \\
    课堂教学技能 & 4.272 & 0.015 & 二级$>$一级* \\
    教学评价技能 & 1.671 & 0.191 & 无显著差异 \\
    教学研究技能 & 0.234 & 0.792 & 无显著差异 \\
    数学教学技能 & 1.462 & 0.235 & 无显著差异
  \end{tblr}
  *在0.05的显著性水平上平均数差异显著。
\end{table}


从表~\ref{table:各维度在计算机等级上的差异性分析结果} 可以看出：（1）计算机等级不同的师范生在数学教学技能上不存在显著的统计学差异，说明计算机等级对师范生的数学教学技能影响不明显。（2）单看每一项教学技能，可以发现计算机等级不同的师范生在课堂教学技能上的差异显著性检验为0.015，明显小于0.05，说明计算机等级不同的师范生在课堂教学技能上存在差异，且计算机等级为二级的师范生平均分明显高于计算机等级为一级的师范生。而计算机等级不同的师范生在其他教学技能上不存在显著的差异，因为sig均大于0.05。



\section{相关性分析}


\subsection{数学教学技能各维度相关性}

本研究对师范生的数学教学技能的五个维度进行了相关性分析，相关系数如表~\ref{table:数学教学技能各维度相关性} 所示。

\begin{table}
  \centering
  \caption{数学教学技能各维度相关性}
  \label{table:数学教学技能各维度相关性}
  \begin{tblr}{
    width  = \textwidth,
    colspec = {
      X[2,c]X[1,c]X[1,c]X[1,c]X[1,c]X[1,c]X[1,c]
    },
    hline{1,2,Z} = {solid},
  }
    & 1       & 2       & 3       & 4       & 5       & 6 \\
    1.教学设计技能 & 1       &         &         &         &         &   \\
    2.数学说课技能 & 0.569** & 1       &         &         &         &   \\
    3.课堂教学技能 & 0.549** & 0.500** & 1       &         &         &   \\
    4.教学评价技能 & 0.563** & 0.514** & 0.582** & 1       &         &   \\
    5.教学研究技能 & 0.496** & 0.281** & 0.362** & 0.522** & 1       &   \\
    6.数学教学技能 & 0.816** & 0.743** & 0.744** & 0.835** & 0.713** & 1 \\
  \end{tblr}
  ** 在 0.01 级别（双尾），相关性显著。
\end{table}


结果显示：数学教学技能的各维度之间均存在显著的正相关，且都有显著性$P=0.000<0.01$。

数学教学技能与各项教学技能的相关性系数分别为0.816、0.743、0.744、0.835、0.713，说明数学教学技能与教学评价技能的相关性最大，其次是教学设计技能，然后是数学说课技能和课堂教学技能，最后是教学研究技能。

教学设计技能与各项教学技能的相关性系数分别为0.569、0.549、0.563、0.496，说明且教学设计技能强的师范生其他教学技能也强，且教学设计技能与数学说课技能的相关性最大，与教学研究技能相关性最小。

数学说课技能与各项教学技能的相关系数分别为0.569、0.500、0.514、0.281，说明数学说课技能强的师范生其他技能也强，且数学说课技能与教学研究技能相关性最小。

课堂教学技能与各项教学技能的相关系数分别为0.549、0.500、0.582、0.362，说明课堂教学技能强的师范生其他技能也强。并且课堂教学技能与教学研究技能相关性最小。

教学评价技能与各项教学技能的相关系数分别为0.563、0.514、0.582、0.522，说明教学评价技能强的师范生其他技能也强。并且教学评价技能与课堂教学技能相关性最大，教学评价技能与数学说课技能相关性最小。

教学研究技能与各项教学技能的相关系数分别为0.496、0.218、0.362、0.522，说明教学研究技能强的师范生其他技能也强，且教学研究技能强与教学评价技能相关性最大。

从整体上来看，课堂教学技能与教学评价技能相关性最大，数学说课技能与教学研究技能相关性最小。


\subsection{课堂教学技能各维度相关性}

本研究对师范生的课堂教学技能的七个维度进行了相关性分析，相关系数如表~\ref{table:课堂教学技能各维度相关性分析} 所示。

\begin{table}
  \centering
  \caption{课堂教学技能各维度相关性分析}
  \label{table:课堂教学技能各维度相关性分析}
  \begin{tblr}{
    width  = \textwidth,
    colspec = {
      X[2,c]X[1,c]X[1,c]X[1,c]X[1,c]X[1,c]X[1,c]
    },
    hline{1,2,Z} = {solid},
  }
    & 1       & 2       & 3       & 4       & 5       & 6       & 7 \\
    1.语言技能   & 1       &         &         &         &         &         &   \\
    2.板书技能   & 0.423** & 1       &         &         &         &         &   \\
    3.讲解技能   & 0.449** & 0.386** & 1       &         &         &         &   \\
    4.提问技能   & 0.436** & 0.435** & 0.509** & 1       &         &         &   \\
    5.导入技能   & 0.553** & 0.462** & 0.500** & 0.568** & 1       &         &   \\
    6.结束技能   & 0.432** & 0.395** & 0.380** & 0.544** & 0.639** & 1       &   \\
    7.媒体使用技能 & 0.402** & 0.479** & 0.373** & 0.426** & 0.540** & 0.455** & 1 \\
  \end{tblr}
  ** 在 0.01 级别（双尾），相关性显著。
\end{table}


结果显示：课堂教学技能的七个维度之间均存在显著的正相关，且都有显著性$P=0.000<0.01$。

语言技能与其他课堂教学技能的相关性系数分别为0.423、0.449、0.436、0.553、0.432、0.402，说明语言技能强的师范生其他课堂教学技能也强，且语言技能与导入技能相关性最大，与媒体使用技能相关性最小。

板书技能与其他课堂教学技能的相关性系数分别为0.423、0.386、0.435、0.462、0.395、0.479，说明板书技能强的师范生其他课堂教学技能也强，板书技能与媒体使用技能相关性最大，与讲解技能相关性最小。

讲解技能与其他课堂教学技能的相关性系数分别为0.449、0.386、0.509、0.500、0.380、0.373，说明讲解技能强的师范生其他课堂教学技能也强，讲解技能与提问技能相关性最大，与媒体使用技能相关性最小。

提问技能与其他课堂教学技能的相关性系数分别为0.436、0.435、0.509、0.568、0.544、0.426，说明提问技能强的师范生其他课堂教学技能也强，提问技能与导入技能相关性最大，与媒体使用技能相关性最小。

导入技能与其他课堂教学技能的相关性系数分别为0.553、0.462、0.500、0.568、0.639、0.540，说明导入技能强的师范生其他课堂教学技能也强，导入技能与结束技能相关性最大，与板书技能相关性最小。

结束技能与其他课堂教学技能的相关性系数分别为0.432、0.395、0.380、0.544、0.639、0.455，说明结束技能强的师范生其他课堂教学技能也强，结束技能与导入技能相关性最大，与讲解技能相关性最小。

媒体使用技能与其他课堂教学技能的相关性系数分别为0.402、0.479、0.373、0.426、0.540、0.455，说明媒体使用技能强的师范生其他课堂教学技能也强，媒体使用技能与导入技能相关性最大，与讲解技能相关性最小。

从整体上来说，结束技能和导入技能的相关性最高，讲解技能与媒体使用技能的相关性最低。